\chapter[Audio analysis \textsuperscript{Kugler}]{Audio analysis}
\label{chap:Audio analysis}

A big part of the project is to give players the ability to play with their own instruments. To archive this we had to face a number of problems. 
\begin{itemize}
	\item Bass guitars can play very low notes which frequencies are closer together and therefore are hard to distinguish.
	\item Guitars have many overtones which could be misinterpreted as notes. 
	\item Microphones and sound \glspl{pick-up} vary a lot in sensitivity. It is impossible to find a threshold that works for everybody. 
	\item Microphones have different sensitivity for different frequencies. The input has to be normalized with an equalizer. 
\end{itemize}

To work well for many people all parameters have to be adjusted at runtime dynamically. Therefore the program generates statistics about the characteristics of the audio signal and calculates the parameters for note detection accordingly. 


The algorithm presented here is similar to the on described in \cite{harmo1} in that way that we use what they call \emph{harmonic atoms}. However, for performance reasons we use the generalised Hough transform instead of a matching pursuit to analyse the frequency spectrum. 


To detect notes we use a method called short-time Fourier transformation. The audio signal is processed in the following steps:
\begin{enumerate}
	\item Split up the signal in multiple overlapping slices.
	\item Apply the Hann window function to each slice.
	\item Generate a frequency spectrum for every slice.
	\item Remove noise and apply an equalizer to account for differences in microphones and \glspl{pick-up}.
	\item Apply a scoring function to remove overtones and increase frequency resolution.
\end{enumerate}

Each step is described in further detail in the following sections.

\section{Choice of audio input parameters}
An audio input basically has two parameters. The sample rate, which is the number of times the signal is measured and the bit rate, the accuracy at which signals are measured.
For example, standard audio CDs use a sample rate of 44.1kHz and a bit rate of 16 bit. 
Since higher sample rate leads to a better detection quality the program tries to use the highest option that is available on the system.

\section{Splitting up the input into slices}
In a first step the continuous audio signal is split up into slices, each containing 4096 samples. The more samples are used the more accurate can the frequency be determined. The range in which frequencies can be found can be calculated as
$$ \frac{samplerate}{N} \leq F \leq \frac{samplerate}{2} $$
where $N$ is the number of samples in the slice and $samplerate$ is the number of samples per second of the audio signal. The upper bound results from the Nyquist–Shannon sampling theorem. Because the \gls{FFT} is a divide and conquer algorithm that repeatedly splits up the samples into halves, the number of samples in the buffer have to be a power of 2.

We chose to use a size of 4096 samples. With the common sample-rate of 44.1 kHz this is high enough to accurately detect the lowest notes of bass guitars. 4096 samples are about 0.1 second at this sample-rate. To get a higher time-resolution, the slices overlap each in other in such a way that a new slice begins every 5 milliseconds. Therefore, the samples are written into a ring-buffer which is copied every time enough new samples have been written (see figure \ref{fig:ring-buffer}).

\begin{figure}
	\includegraphics[width=\textwidth]{ring_buffer}
	\caption{An example of the content of the ring-buffer while a bass guitar is playing. Note the point of discontinuity where the new samples are written over the old ones. (Information about the scale of the Y axis is not available as we do not know the sensitivity of the microphone)}
	\label{fig:ring-buffer}
\end{figure}

\section{Hann window function}
The \gls{FFT} interprets the buffer as a repeating signal. To avoid noise that is generated by the point of discontinuity we apply a window function that is zero-valued at that point. 
We chose the Hann function for this purpose because it has very low aliasing while only slightly decreasing resolution. It is defined as 
$$ w(n)=  \sin^2 \left ( \frac{ \pi n}{N-1} \right) $$
where $n$ is the current samples index and $N$ the number of samples in the slice.
To apply it to the signal we multiply each sample with the corresponding value of the Hann function.

After applying the Hann window function and shifting along the x axis to the example buffer above, we get the buffer that will be used for the FFT (see figure \ref{fig:Hann window}).
\begin{figure}
	\includegraphics[width=\textwidth]{windowed_buffer}
	\caption{What the signal looks like after applying the Hann window function. The buffer has been rotated so that the point of discontinuity is at $x=0$. Since the Hann window function is 0 at the first and the last point we get a continuous signal suitable for the frequency analysis.}
	\label{fig:Hann window}
\end{figure}

\section{Frequency analysis}
\index{Fast Fourier Transformation}
The next step of the audio processing is the \gls{FFT}. 
The \gls{FFT} is a method to transform a signal from time domain to frequency domain. The result is a list of complex numbers that represent the amplitude and phase of each frequency. 
We used the {\tt KissFFT} library (see \ref{sec:KissFFT}) for the calculation. 
Because the \gls{FFT} is meant to be used on stationary signals while we are dealing with a continuous signal, we split up the input into slices, each a few milliseconds in length and apply a window function to them. 
Samples from the audio device are written into a ring buffer and once a certain amount of samples have been written the \gls{FFT} is applied to them. The frequency spectra overlap each other to get a higher time resolution while maintaining a high frequency resolution. 

\section{Removing noise}
Most noise in audio signals appears as an always prevalent amplitude on certain frequencies. To remove it, the minimum amplitude that has been measured on a frequency so far is subtracted from the measurement. 

To achieve this, the program keeps a list with the minimum amplitudes for multiple frequency ranges.

\section{Automatic equalizer}
Additionally to the difference in overall sensitivity, cheap pick-ups and microphones also show different sensitivity for different frequencies. For better note detection an equalizer is applied to the signal. 

An equalizer amplifies or tones down individual frequencies. Since the spectrum of guitars spread rather evenly over the whole frequency range the sensitivity can be measured by finding the maximum amplitude that is measured on a frequency range so far. 

Additionally to the list of minima that is used for the noise removal, the program also keeps a list of the maximum amplitude for multiple frequency ranges. The equalized amplitude for a frequency is calculated as 
$$a'(F, t) = \frac{a(F, t) - min(F)}{max(F) - min(F)}$$
where $a(F, t)$ is the amplitude of frequency $F$ at time $t$ and $min(F)$ and $max(F)$ are the minimum and maximum amplitudes measured in the range $F$ lies in.

\section{Applying the scoring function}
Every natural sound consists of a fundamental frequency and a set of overtones whose frequencies are decimal multiples of the fundamental frequency (see figure \ref{fig:Bass spectrum}). Since the overtones can (and the first one always will) fall on the frequencies of other pitch values they have to be filtered from the signal. We use a scoring function to determine how likely a given frequency is a note.

\begin{figure}
	\includegraphics[width=\textwidth]{bass_spectrum}
	\caption{The frequency spectrum of a bass guitar playing the note G3. Notice the fundamental frequency at about 200 Hz and the overtones at multiples of that frequency.}
	\label{fig:Bass spectrum}
\end{figure}

To remove overtones and also increase frequency resolution we apply a scoring function that is defined in such a way that frequencies with overtones but no undertones get the highest score. 
The score for a pitch is defined as
$$ score(x) = \sum_{i = 1}^{I}\frac{F(x * i)}{\sqrt{i}} - \sum_{i = 2}^{I}\frac{F(x / i)}{\sqrt{i}} $$
where $x$ is the current frequency, $F(x)$ is the amplitude of frequency $x$ and $I$ is the number of overtones to look for.

Possible notes will show up as local maxima in the resulting function.
Because of accuracy issues and noise it often also detects notes near the actual note. To account for this, we only consider frequencies that are a maximum  within a certain radius. For a radius we chose one semi-tone because accords with notes that close together (a {\em major second}) are very unusual in modern music. 

\section{Automatic thresholds}
After potential pitches have been determined, they are compared to a threshold. The optimal value for this varies between different instruments and audio hardware. Multiple algorithms to find the best threshold have been tried out. These are described here.

\subsection{Finding the two most distant clusters}
The most basic algorithm to find a threshold is implemented in the class {\tt AutoThreshold}. Basically, it consists of a single function {\tt bool is(int value)}, which compares a value to the threshold and adjusts the threshold. 

For reasons of efficiency, the threshold is not updated every time a new value is compared but only after a call to the function {\tt update()}.

The algorithm underlies the idea that the distribution of measurements creates two cluster points. By counting how often each value has been compared it tries to find these clusters and then determines the threshold. 

To find the optimal threshold the algorithm splits up the set of measured amplitudes in such a way that the emerging subsets' averages are as far apart as possible. This is achieved by iterating over all possible thresholds.

The ideal threshold is the one for which the following expression reaches its maximum 
$$ \frac{\sum_{i = t}^{N - 1}(i * P(i))}{N - t} - \frac{\sum_{i = 0}^{t - 1}(i * P(i))}{t} $$
where $t$ is the index of the first range that lies just above the threshold, $N$ is the number of ranges and $P(i)$ is the number of measurements that fall into the range at index $i$.

Since the ratio between silence and playing can vary the algorithm should still work if one of the two clusters is much higher. This is where this method fails. If the clusters are not of similar size the distance between them can be increased by moving the threshold beyond the smaller.

\subsection{Finding the two narrowest clusters}
This algorithm underlies the same principles the previous one does but instead of trying to have the two clusters as distant as possible they should be as narrow as possible. Therefore we calculate the variance of both subsets and try to minimize their sum.

First we calculate the average values of both sets with 
$$ \mu_{0} = \frac{\sum_{i = 0}^{t - 1}i * P(i)}{t} $$
$$ \mu_{1} = \frac{\sum_{i = t}^{N - 1}i * P(i)}{N - t} $$
and then the sum of the variances of both subsets as
$$ \sum_{i = 0}^{t - 1}P(i) * (i - \mu_{0})^{2} + \sum_{i = t}^{N - 1}P(i) * (i - \mu_{1})^{2} $$
where $t$ is the index of the first range that lies just above the threshold, $N$ is the number of ranges and $P(i)$ is the number of measurements that fell into the range at index $i$.

This approach is only useful when the notes have similar volume since else the cluster for the upper cluster spread to much. Additionally, this solution is rather slow even on modern computers.

\subsection{Considering the number of notes we are looking for}
This is a very different approach as it does not rely on statistics about the signal but instead considers how many notes an instrument can play at once. Most guitars have 4 to 6 strings and therefore can play only that many notes at once. Additionally we assume that all simultaneous notes have about the same volume. 

After the maxima in the scoring function have been found we look for the highest score and consider every other maxima that is not much quieter than that as a note. If the number of notes found is bigger than 6 we consider the signal noise as it can not come from a guitar.
